% Copyright (c) 2026 Eugene Joseph M. Ragasa.
% SPDX-License-Identifier: Apache-2.0
% Author: Eugene Joseph M. Ragasa
% Formal proof completed: 2026-10-30
\documentclass[11pt]{article}
\usepackage[T1]{fontenc}
\usepackage{amsmath,amssymb,booktabs,tabularx}
\usepackage[margin=0.85in]{geometry}
\usepackage[hidelinks]{hyperref}
\newcolumntype{Y}{>{\raggedright\arraybackslash}X}
\title{Proof Dossier: Four-Stage Root-Mean-Square Bound}
\author{Eugene Joseph M. Ragasa}
\date{Formal proof completed: 2026-10-30}
\begin{document}
\raggedright
\maketitle
\section{Identity and status}
\begin{tabularx}{\textwidth}{@{}lY@{}}
\toprule
Statement ID & \texttt{BG-HILBERT-THM-025} \\
Status & proved; no admissions \\
Lean declaration & \path{ProjectKoios.BandGauge.Hilbert.finiteRootMeanSquareNorm_totalResidual_le_three} \\
Owning module & \path{ProjectKoios/BandGauge/Hilbert/FourStageResidualChannels.lean} \\
Author and license & Eugene Joseph M. Ragasa; \texttt{Apache-2.0} \\
\bottomrule
\end{tabularx}

\section{Exact mathematical statement}
For the four-stage residual fields over a finite inventory $I$, with values in a normed additive commutative group, let
\[
 \operatorname{RMS}(r)=\frac{\bigl(\sum_{i\in I}\|r_i\|^2\bigr)^{1/2}}{\sqrt{|I|}}.
\]
Using totalized division when $I$ is empty,
\[
 \operatorname{RMS}(r_{\rm total})\le
 \operatorname{RMS}(r_{\rm int})+
 \operatorname{RMS}(r_{\rm conv})+
 \operatorname{RMS}(r_{\rm trunc}).
\]

\section{Immutable Formal provenance}
The declaration is developed against Formal base commit
\texttt{0062d687c1306c1df85d9c674fdf0239d05a8624}, tree
\texttt{cd60a1bb5cd005547338398621be3c596c865aa4}. The owning-module SHA-256 is
\texttt{3931c865c24db66dd42ac65c922742448402134dbc071e596df686d124e3dfec}; the locator is
\texttt{FourStageResidualChannels.lean:107} and \texttt{BG-HILBERT-THM-025}. A topic commit is not claimed before checkpoint authorization.

\section{Assumptions, conventions, and statement delta}
The finite inventory and norm are abstract. The denominator is the square root of the inventory cardinality. This extends the earlier two-component RMS triangle bound by retaining conversion and deliberate truncation as separate component fields. It asserts an inequality, not scalar additivity.

\section{Proof structure}
The exact field decomposition gives $r_{\rm total}=r_{\rm int}+r_{\rm conv}+r_{\rm trunc}$. These fields inhabit the finite $\ell^2$ function space. Applying its norm triangle inequality twice bounds the numerator by the sum of the three component numerators. Division by the nonnegative common factor $\sqrt{|I|}$ preserves the inequality, including the totalized empty case.

\section{Lean realization and trusted base}
Lean represents the finite $\ell^2$ space with \texttt{PiLp 2}, proves the field equality extensionally, and applies two norm triangle inequalities before a nonnegative-division lemma. Toolchain: Lean \texttt{4.35.0} at \texttt{eb6cafb7290e8bd5f2937522d8d804f58f02bafc}; Mathlib \texttt{56a005855ea17e0b0958f73ec801ae140b0a1679}.
\begin{verbatim}
#print axioms ProjectKoios.BandGauge.Hilbert.
  finiteRootMeanSquareNorm_totalResidual_le_three
-- [propext, Classical.choice, Quot.sound]
\end{verbatim}
No repository-owned admission occurs.

\section{Scientific, numerical, and certificate boundaries}
The theorem proves no RMS value for an external dataset and authenticates no mesh cardinality, runtime array, binary64 norm, or stage implementation. It supplies no interval enclosure or proposition-specific certificate. The current checked runtime-certificate count remains zero.

\section{AI, automation, and review record}
The proof and dossier were AI-assisted and checked by the maintainer. Lean proof checking is not independent scientific, mathematical, novelty, bibliography, or publication review.

\section{Review and revision history}
\begin{tabularx}{\textwidth}{@{}lllY@{}}
\toprule Date & Revision & Process & Outcome and limitation \\
\midrule 2026-10-30 & initial dossier & AI-assisted maintainer review & RMS normalization and inequality boundary checked; no runtime promotion \\
\bottomrule
\end{tabularx}
\end{document}
